\documentclass[a4paper,11pt]{article}

\usepackage[margin=2.5cm]{geometry}
\usepackage[T1]{fontenc}
\usepackage[utf8]{inputenc}
\usepackage{lmodern}
\usepackage{booktabs}
\usepackage{tabularx}
\usepackage{longtable}
\usepackage{array}
\usepackage{enumitem}
\usepackage{hyperref}

\hypersetup{
  colorlinks=true,
  linkcolor=black,
  urlcolor=blue,
  citecolor=black
}

\setlength{\parindent}{0pt}
\setlength{\parskip}{0.45em}
\renewcommand{\arraystretch}{1.15}
\newcolumntype{Y}{>{\raggedright\arraybackslash}X}

\title{\textbf{Benchmarking LLM-Based Honeypots with a Scripted Local Harness:
Setup, Findings and Limitations (v2)}}
\author{Muhammad Rumman Aslam}
\date{}

% Table sources (kept in sync with each caption):
% tab:targets   harness/targets.yaml; results/exp/stage4/containers.json
% tab:config    docker/*.compose.yaml; BENCHMARK.md; docs/PATCH_NOTES.md
% tab:v1-ssh    results/summary.md; results/exp/summary_v2.md
% tab:floor     results/exp/floor.md
% tab:abac      results/exp/validation/{target}_comparison.csv
% tab:s4-ssh    results/exp/stage4/{latency,determinism,bootstrap_ci,extra_vs_gt}.csv
% tab:s4-vs-v1  results/summary.md; results/exp/summary_v2.md; results/exp/stage4/latency.csv
% tab:galah     results/exp/stage4/analysis.md; analysis_followup.md S5; results/summary.md
% tab:quiet     results/exp/stage4/quiet_check_full.csv (full 288 rows)
% tab:timeouts  results/exp/stage4/analysis_timeouts.md
% tab:maxtokens results/exp/stage4/analysis_timeouts.md Part A.2; BENCHMARK.md
% tab:lifetime  results/exp/lifetime_control/beelzebub_persistent/lifetime.json

\begin{document}
\maketitle

\begin{center}
\textbf{Baseline batch: Nov 2026, 5 runs per target, 120\,s timeout, quiet
timing, one long-lived container per target. Main results: Stage 4, Nov 2026,
20 runs per target, prompt timing, fresh container per slot, shuffled order.
Version 2 adds the 20-run Stage 4 pass, timeout attribution against the API
proxy log, a per-target max-token check, and a single-container
Beelzebub lifetime control.}
\end{center}

\begin{abstract}
This report documents a scripted local benchmark of LLM-based SSH and HTTP
honeypots, a static Cowrie baseline, ground-truth services, an ANSI-hidden
prompt-injection detector target, and HoneyBee as a separate environment
generator. It reports two batches. The baseline batch (v1) used five runs per
target in quiet timing against one long-lived container per target. The main
results come from Stage 4: twenty runs per target in prompt timing, a shuffled
target order fixed by seed 2026, a fresh container for every slot, write-time
credential masking, and a local API proxy that logged 3,699 requests at a
corrected cost of \$2.02. Each SSH run sent 37 ordered commands, of which 36
were scored; each HTTP run sent 20 requests; the runner timeout was 120\,s and
timed-out observations are retained as censored events. Similarity comparisons
use ANSI-stripped prefixes capped at 4,000 characters. In Stage 4 the LLM SSH
targets had median latencies of 2.18--2.39\,s against 0.015\,s for the
ground-truth service (values below 50\,ms are limited by the 15.6\,ms Windows
clock tick), 0--45 timeouts out of 720 commands, masked determinism of
0.509--0.848 against 0.992 for ground truth, and masked fidelity of
0.433--0.637. Timeout attribution shows that 9 of the 86 timed-out rows were
API-side (241--291\,s, all ending at exactly 65,536 completion tokens), 66 were
target-side (API duration 1.5--28.1\,s, median 5.4\,s; for the 51 VelLMes rows
4.5--7.8\,s with exactly 800 completion tokens; the reply never reached the
harness), and 11 had no matching request. Observed completion tokens respect
the documented caps (4,096 for AdvancedshelLM, 800 for VelLMes); the uncapped
targets reach the model ceiling of 65,536. A Beelzebub container-lifetime
control (ten sessions in one container) shows a rising median (3.88 to
5.95\,s, Spearman $\rho=0.915$) and 6/360 timeouts, all in the three sessions
that contained a 65,536-token generation. These are scripted measurements from
one shared model and endpoint; they do not rank tools, establish causes, or
generalise beyond the tested configurations.
\end{abstract}

\section{Introduction and goals}

LLM-based honeypots can answer shell or HTTP interactions with outputs
constructed by a language model rather than by a fixed emulation layer. A
benchmark must therefore record both conventional service observations and
properties specific to generated text: repeat consistency, output length,
interior prompt leakage, empty or unexpected outputs, and timeout behaviour.
This project uses a local harness to make those observations comparable across
several SSH and HTTP targets.

\paragraph{Related work.} Vetterl and Clayton fingerprint low- and
medium-interaction honeypots at internet scale~\cite{bitterharvest}. Sladi\'c
et al.\ describe generative honeypots built on LLMs~\cite{llmintheshell} and
the VelLMes high-interaction deception framework~\cite{vellmes-paper}. Wang et
al.\ present HoneyGPT for terminal honeypots~\cite{honeygpt}. Reworr and Volkov
report a honeypot that monitors AI hacking agents in the wild~\cite{reworr},
and Bridges et al.\ survey honeypots and LLMs~\cite{bridgessok}. This report
does not test their claims; it measures the local configurations below.

The study is small in scale: 20 runs per target after 5 in the published batch,
one shared model and API endpoint for the LLM targets, a fixed command/request
corpus, and two batches whose protocols differ. The goal is measurement and reproducibility rather than a
universal ranking. The ground-truth SSH service is Ubuntu 20.04 with OpenSSH,
and the HTTP reference is Ubuntu 20.04 with nginx. Cowrie supplies a static SSH
baseline. The llm-honeypot target is treated as a detector target and excluded
from fidelity, similarity, determinism, and aggregate averages.

Three practical goals shape the report: document the local setup and target
configuration; state the metrics and accounting rules precisely, including the
differences between the two batches; and preserve the limitations and
unresolved questions that should guide a larger study. Every number below is
traced to a file under \texttt{results/}; items the data cannot settle are
marked UNVERIFIED.

\section{Setup}

\subsection{Targets and ports}

All benchmark services were published on loopback ports. Table~\ref{tab:targets}
summarises the target inventory. HoneyBee is deliberately listed as a
generator: it produces Docker-oriented artifacts and is not an interactive
honeypot in this benchmark. The implementations are DECEIVE~\cite{deceive},
Beelzebub~\cite{beelzebub}, Cowrie~\cite{cowrie}, shelLM/AdvancedshelLM~
\cite{advancedshellm}, VelLMes~\cite{vellmes} and galah~\cite{galah}.

\begin{table}[ht]
\centering
\caption{Final target inventory. Source: \texttt{harness/targets.yaml};
container identities in \texttt{results/exp/stage4/containers.json}.}
\label{tab:targets}
\small
\begin{tabularx}{\textwidth}{l l r l Y}
\toprule
Target & Track & Port & Role & Recorded scope \\
\midrule
DECEIVE & SSH & 8022 & LLM honeypot & Scripted SSH responses \\
Beelzebub & SSH & 2222 & LLM honeypot & Scripted SSH responses \\
Cowrie & SSH & 2223 & Static baseline & Static emulation \\
AdvancedshelLM & SSH & 2224 & LLM honeypot & Subprocess/wrapper path \\
Ground truth SSH & SSH & 2225 & Reference & Ubuntu 20.04/OpenSSH \\
llm-honeypot & SSH & 2227 & Detector target & ANSI-hidden marker detector \\
VelLMes & SSH & 2228 & LLM honeypot & Scripted SSH responses \\
galah & HTTP & 8088 & LLM HTTP responder & Scripted HTTP responses \\
Ground truth HTTP & HTTP & 2226 & Reference & Ubuntu 20.04/nginx \\
HoneyBee & generator & 2229 config & LLM generator & Output-quality evaluation \\
\bottomrule
\end{tabularx}
\end{table}

\subsection{Runtime and configuration}

The documented reproduction environment is Windows 10/11 with Python 3.10+
(tested with Python 3.12), Git, and Docker Desktop with WSL 2. Services are
bound to loopback. The LLM targets use one shared model/API endpoint, while
the configuration values that were explicitly recorded differ by target.
Table~\ref{tab:config} preserves ``not recorded'' where the permitted
documents did not specify a value. Sources: the compose files under
\texttt{docker/}, \texttt{BENCHMARK.md} (target configuration table) and
\texttt{docs/PATCH\_NOTES.md}.

{\footnotesize
\setlength{\tabcolsep}{3pt}
\begin{longtable}{@{}>{\raggedright\arraybackslash}p{2.45cm}
                    >{\raggedright\arraybackslash}p{2.95cm}
                    >{\raggedright\arraybackslash}p{2.25cm}
                    >{\raggedright\arraybackslash}p{1.35cm}
                    >{\raggedright\arraybackslash}p{1.35cm}
                    >{\raggedright\arraybackslash}p{4.15cm}@{}}
\caption{Recorded target configuration. Sources: compose files under
\texttt{docker/}, \texttt{BENCHMARK.md}, \texttt{docs/PATCH\_NOTES.md}.}
\label{tab:config}\\
\toprule
Target & Stack & Model & Temp. & Max tokens & User/host and main configuration \\
\midrule
\endfirsthead
\toprule
Target & Stack & Model & Temp. & Max tokens & User/host and main configuration \\
\midrule
\endhead
\bottomrule
\endfoot
DECEIVE & Python 3.11-slim + DECEIVE & glm-5p3-flash & 0.2 & no cap &
guest / not recorded; output stripping, Fireworks env, exit sentinel \\
Beelzebub & Go image + mounted YAML & glm-5p3-flash & not exposed & not exposed &
guest / ubuntu; API-key and host fallback, benchmark YAML \\
AdvancedshelLM & Python 3.11-slim + shelLM & glm-5p3-flash & 0.2 & 4096 &
guest / not recorded; client, exit handling, env fallback, wrapper \\
VelLMes & Python 3.11-slim + VelLMes & glm-5p3-flash & 0.2 & 800 &
user / not recorded; SSH wrapper and benchmark config \\
galah & Go 1.22 build + Debian slim & glm-5p3-flash & 0.2 & not recorded &
not applicable / not recorded; response cleaning, entrypoint \\
Cowrie & official image digest & not applicable & not applicable & not applicable &
root / ubuntu; mounted config, no upstream source edit \\
llm-honeypot & Cowrie v2.6.1 overlay & not applicable & not applicable & not applicable &
root / not recorded; overlay Dockerfile and v2.6.1 pin \\
Ground truth SSH & Ubuntu 20.04 + OpenSSH & not applicable & not applicable & not applicable &
guest / ubuntu; PTY and PAM changes \\
Ground truth HTTP & Ubuntu 20.04 + nginx 1.18.0 & not applicable & not applicable & not applicable &
not applicable / not recorded; Dockerfile and compose \\
HoneyBee & Python 3.11 + Streamlit & glm-5p3-flash & 0.2 & 8192 &
not applicable; deploy path disabled \\
\end{longtable}
}

Third-party source changes are documented as patches and notes. The verified
third-party commit hashes are the HoneyBee clone
\texttt{dbf62c42604eeae2cfbfd4d36b99bef65555c303} and the VelLMes clone
\texttt{192fa5084df84b38a16f6581939fd3143c350181}. \textbf{The DECEIVE clone
commit is UNVERIFIED (no \texttt{.git} in the local clone).} There is no
\texttt{.git} directory in the local \texttt{DECEIVE/} working tree, so a
\texttt{rev-parse} inside it resolves to this benchmark repository and cannot
read a DECEIVE commit.

The upstream repository cited in v1 was
\texttt{https://github.com/coverflex-tech/DECEIVE}, which does not resolve
(``Repository not found'', checked 21 November 2026); the upstream used here
is \texttt{https://github.com/splunk/DECEIVE}. A fresh clone of
\texttt{splunk/DECEIVE} on the same date has HEAD
\texttt{cce98f3dcd9de74f32f683998e95aa19b795724e} (2026-05-13, MIT licence).
No upstream revision of \texttt{SSH/ssh\_server.py} is byte-identical to the
local file: the closest is that HEAD revision (\texttt{1351d54}, 75 changed
lines, matching the ten hunks recorded in \texttt{docs/PATCH\_NOTES.md}~S1),
and the two \texttt{LICENSE} files are identical in text. The upstream
template is \texttt{SSH/config.ini.TEMPLATE}; the local \texttt{SSH/config.ini}
differs from it only in \texttt{model\_name}. Because no revision matches byte
for byte and the clone kept no \texttt{.git}, the exact clone commit stays
UNVERIFIED. Other clone commit hashes are omitted here. The ground-truth
services are benchmark-created services, not third-party source repositories.

HoneyBee is not a honeypot in this study. It is a Streamlit application that
uses an LLM to generate misconfigured Docker artifacts. The one-click deploy
path is disabled, the image has no Docker CLI or Docker socket mount, and the
evaluation reports generator validity rather than a honeypot security
assessment.

The harness, the documentation and this report are MIT-licensed (see
\texttt{LICENSE}); each third-party tool keeps its own licence, listed in
\texttt{THIRD\_PARTY\_NOTICES.md}, and the raw data collected from the
third-party honeypots are published for reproducibility.

\subsection{DECEIVE configuration evidence}

Four configuration items were checked in the local clone, and the two
\texttt{config.ini} files in the repository were compared:
\begin{itemize}[leftmargin=*]
\item \texttt{DECEIVE/SSH/config.ini} line 32: \texttt{model\_name =
  accounts/fireworks/models/glm-5p3-flash}; line 61: \texttt{temperature =
  0.2}; line 67: \texttt{trimmer\_max\_tokens = 64000}.
\item \texttt{recovered/deceive\_app/SSH/config.ini}, a copy taken from the
  running DECEIVE application, carries the same \texttt{model\_name} (line 32)
  and \texttt{temperature = 0.2} (line 72), and additionally sets
  \texttt{tiktoken\_model\_name = gpt-4o} (line 43). The working clone omits
  that key and relies on the code default: \texttt{ssh\_server.py} line 420
  reads \texttt{config['llm'].get('tiktoken\_model\_name', 'gpt-4o')}, so both
  files resolve to the same value in the API request.
\item \texttt{DECEIVE/SSH/ssh\_server.py} line 429: \texttt{ChatOpenAI(**kwargs)}
  is built from \texttt{model}, \texttt{temperature}, \texttt{api\_key} and
  \texttt{tiktoken\_model\_name} (lines 421--428) with no \texttt{max\_tokens}
  argument, so the model receives no output-length cap per request. This is
  consistent with the oversized outputs observed in both batches
  (Section~\ref{sec:s4-quality}).
\item The clone is patched for this benchmark, as recorded in
  \texttt{docs/PATCH\_NOTES.md}~S1: output stripping, \texttt{OPENAI\_API\_KEY}
  and \texttt{OPENAI\_BASE\_URL} resolution, the \texttt{tiktoken\_model\_name}
  key, exit-sentinel matching and lazy provider imports. The code reads the
  model name only from \texttt{config.ini}; no environment variable overrides
  it.
\end{itemize}

Which of the two \texttt{config.ini} files is byte-identical to the file that
ran either batch is \textbf{UNVERIFIED}. Both name the same model and the same
temperature and differ only in comments and in the
\texttt{tiktoken\_model\_name} key, whose effective value is identical; no
hashed copy of the application directory was retained for either batch.

\section{Baseline batch (v1): protocol and results}

\subsection{Protocol}

The SSH manifest contains 37 ordered commands: 36 scored commands plus
\texttt{exit}. Its observed categories are recon, files, process, errors,
permissions, and state. Host-specific commands are excluded from fidelity
scoring. The HTTP manifest contains 20 requests: 17 GET, one HEAD,
one OPTIONS, and one POST, including three repeated \texttt{GET /} requests.
Fidelity and determinism comparisons use ANSI-stripped output prefixes capped
at 4,000 characters (\texttt{harness/analyze.py}).


The baseline batch collected five runs per SSH and HTTP target in the
harness's default \emph{quiet} timing mode, with a 120-second runner timeout
and a duplicate guard that skipped complete target/run pairs. Containers were
started once with \texttt{scripts/up\_all.ps1} before the batch and were not
restarted between runs, so all five runs of a target used one long-lived
container (\texttt{docs/STAGE2\_NOTES.md}, published-batch section;
\texttt{results/exp/stage4/analysis\_followup.md}~S3). Raw-file timestamps show
a two-pass SSH order per run and galah collected after the first SSH pass.

\subsection{SSH results}

The scored SSH denominator is 180 commands per target.
Table~\ref{tab:v1-ssh} reports completed observations separately from
timeouts.

\begin{table}[ht]
\centering
\caption{Baseline (v1) SSH results: 5 runs, quiet timing, 180 scored commands
per target. Source: \texttt{results/summary.md}; determinism columns from
\texttt{results/exp/summary\_v2.md}.}
\label{tab:v1-ssh}
\scriptsize
\begin{tabularx}{\textwidth}{l r r r r l r l}
\toprule
Target & Median/p95 & Comple\-ted & Time\-outs & Empty & Fidelity (95\% CI) & Det. & Oversized \\
\midrule
Ground truth & 0.562/0.579 & 180 & 0 & 35 & --- & --- & 0 \\
Cowrie & 0.563/0.610 & 180 & 0 & 40 & 0.5877 (0.4203--0.7438) & 0.9874 & 0 \\
AdvancedshelLM & 1.922/4.504 & 180 & 0 & 26 & 0.5947 (0.4336--0.7386) & 0.7836 & 0 \\
VelLMes & 1.687/4.874 & 171 & 9 & 25 & 0.6070 (0.4601--0.7506) & 0.8364 & 0 \\
DECEIVE & 1.984/10.361 & 175 & 5 & 27 & 0.4731 (0.3493--0.5858) & 0.4664 & 2 \\
Beelzebub & 5.023/31.656 & 164 & 16 & 16 & 0.3737 (0.2319--0.5239) & 0.6800 & 9 \\
llm-honeypot & detector & 180 & 0 & 30 & excluded & excluded & 0 \\
\bottomrule
\end{tabularx}
\end{table}

Timeouts were command-specific. VelLMes \texttt{ps\_aux} timed out in 5/5 runs
and \texttt{env} in 4/5. Beelzebub \texttt{state\_cat\_marker},
\texttt{state\_echo\_marker}, and \texttt{state\_ls\_tmp\_3} each timed out in
3/5 runs. DECEIVE \texttt{cpuinfo} and \texttt{ps\_aux} timed out in runs 1
and 4. The Beelzebub rows total 17 because the table includes the run~1
\texttt{exit} control row; the scored total is 16. DECEIVE totals 5 and VelLMes
totals 9 (\texttt{docs/FACTS\_CHECK.md}).

\subsection{HTTP results}

Galah produced 11/100 5xx responses, six errors including three timeouts, six
distinct \texttt{Server} header values, 12 header-name-set inconsistencies,
0.2364 repeated-\texttt{GET /} similarity, 25/94 status agreements with nginx,
five documentation-example secret hits, and 19 target-port mentions. Its
latency median/p95 was 3.531/8.887 seconds. Ground-truth nginx produced zero
5xx responses, a constant \texttt{Server} header, 1.0000 repeat similarity,
and a 0.015-second median (\texttt{results/summary.md}).

\textbf{Withheld raw bodies.} Nine galah HTTP raw response files from the
baseline batch and the pilot are not published. They contain
model-fabricated credential-like strings (for example Stripe-style
\texttt{sk\_live\_...} values) that GitHub secret scanning flags, and their
provenance cannot be certified: galah's responses are generated text, so the
author cannot show whether a string is fabricated or memorised from training
data. Publishing them would therefore be worse than withholding them, and
rewriting the bodies would break the claim that the raw output is unmodified.
Their paths, sizes and SHA-256 values are recorded in
\texttt{docs/WITHHELD\_FILES.md}; the files are available to reviewers on
request. Because of this, the galah counts of the baseline batch above cannot
be recomputed from the repository alone. The Stage 4 galah files were
collected with write-time credential masking, contain no such strings, and are
published under \texttt{results/exp/stage4/raw\_http/}.

\subsection{HoneyBee results}

HoneyBee was evaluated with \texttt{glm-5p3-flash}, temperature 0.2, maximum
8192 tokens, and three generations per generator/case cell, for 18
generations. Docker Compose validity was 6/6 cases, Dockerfile validity 4/6,
and Nuclei validity 3/6. Length-limited results were 5/18, all in the
PostgreSQL case. Grafana was clean in 9/9 generations and PostgreSQL in 4/9.
The scorer was revised after the first results; original outputs remain in the
recorded validity file. These are generator-validity observations, not a
security verdict.

\section{Stage 2: measurement additions}

Stage 2 added the measurement machinery that Stage 4 then used: an explicit
prompt-timing mode, credential masking, a local API proxy, masked determinism,
and a measured harness latency floor. Details are in
\texttt{docs/STAGE2\_NOTES.md}.

\subsection{Latency floor}
\label{sec:floor}

The harness's default command-completion rule records a command as finished
only after 0.5 seconds of silence following the last received prompt. This
\emph{quiet period} is a property of the harness, not of the server, so every
latency in Table~\ref{tab:v1-ssh} and the 3.0x--8.9x published ratios include a
0.5-second harness floor. The measured floor on the two non-LLM containers is
shown in Table~\ref{tab:floor}; the ground-truth servers themselves reply in
3--9\,ms.

\begin{table}[ht]
\centering
\caption{Measured harness latency floor, 30 trivial commands per cell (from
\texttt{results/exp/floor.md}). Values are quantised by the 15.6\,ms Windows
timer tick.}
\label{tab:floor}
\scriptsize
\begin{tabularx}{\textwidth}{l l r r r r r}
\toprule
Target & Timing mode & n & Min (s) & Median (s) & p95 (s) & Max (s) \\
\midrule
groundtruth\_ssh & quiet & 30 & 0.546 & 0.547 & 0.563 & 0.578 \\
groundtruth\_ssh & prompt & 30 & 0.015 & 0.0158 & 0.0163 & 0.0163 \\
cowrie & quiet & 30 & 0.547 & 0.563 & 0.578 & 0.579 \\
cowrie & prompt & 30 & 0.015 & 0.0155 & 0.0314 & 0.0319 \\
\bottomrule
\end{tabularx}
\end{table}

The same runs in full 37-command form give a quiet median of 0.562\,s and a
prompt median of 0.016\,s on both non-LLM targets
(\texttt{results/exp/floor.md}, \texttt{results/exp/timing\_audit/}). Quiet and
prompt timing are separate code paths (\texttt{\_read\_quiet} versus
\texttt{\_read\_until\_prompt\_settled} in \texttt{harness/run\_ssh.py}), so the
baseline batch is unaffected; quiet remains the default. Stage 4 uses prompt
timing, so its latencies do not contain the 0.5\,s floor.

\subsection{Prompt-mode validation on LLM SSH targets}
\label{sec:prompt-val}

Prompt mode was previously validated only on the two non-LLM targets. Stage 2
ran the full 37-command list in three passes---quiet (A), prompt (B), quiet
again (C)---on all four LLM SSH honeypots with fresh containers per pass, all
via the API proxy. Total API calls: 452 out of a 500-call budget
(\texttt{results/exp/validation/}).

The key diagnostic metric is \texttt{prompt\_without\_echo}: this flag is set
when the harness ends a read on a buffer that contains nothing but a prompt,
the signature of a stale prompt or early cutoff. All four targets returned zero
\texttt{prompt\_without\_echo} rows in pass~B, and no prompt-like line appeared
inside any prompt-mode output. No early cut was observed in 37 commands x 1
pass per target. Stage 4 extends the same check to 5,180 prompt-mode SSH rows,
5,040 of them scored (Section~\ref{sec:s4-quality}).

An A-vs-B comparison against the A-vs-C quiet baseline covers 36 commands per
target. Table~\ref{tab:abac} counts commands whose A-vs-B masked similarity is
below their own A-vs-C value.

\begin{table}[ht]
\centering
\caption{Commands whose A-vs-B masked similarity is below the A-vs-C value for
the same command (from
\texttt{results/exp/validation/\{target\}\_comparison.csv}).}
\label{tab:abac}
\scriptsize
\begin{tabular}{l r r r r}
\toprule
Target & Checked & A-vs-B $<$ A-vs-C & of those, drop $>$ 0.10 & of those, B shorter than A \\
\midrule
DECEIVE        & 36 & 13 & 12 & 7 \\
Beelzebub      & 36 & 14 & 12 & 2 \\
VelLMes        & 36 & 9  & 0  & 0 \\
AdvancedshelLM & 36 & 6  & 2  & 2 \\
\bottomrule
\end{tabular}
\end{table}

Every one of the 26 commands with a drop above 0.10 has \texttt{prompt} set,
\texttt{prompt\_without\_echo} false, no timed-out pass~B, and no prompt-like
line inside its output, so none shows the early-cut signature. Fifteen of the
26 had a longer pass-B output than pass~A. The A-vs-C baseline is a single
quiet-vs-quiet comparison per command and so does not bound LLM variance:
whether the residual differences are measurement artefacts or ordinary model
variation is UNVERIFIED.

\subsection{Masked determinism}

The v1 determinism column compares outputs character by character, so
timestamps, PIDs, uptime values and hostnames contribute. Stage 2 replaces
those fields with placeholders before the comparison. On the baseline batch,
masking moves every target by at most 0.04
(\texttt{results/exp/summary\_v2.md}): ground truth 0.9746 to 0.9841, Cowrie
0.9874 to 0.9957, AdvancedshelLM 0.7836 to 0.8034, VelLMes 0.8364 to 0.8472,
Beelzebub 0.6800 to 0.6815, DECEIVE 0.4664 to 0.5038. Masking volatile fields
did not close the gap between the static baseline and the LLM targets. The
ground-truth service is itself not perfectly deterministic (0.9746 raw), so
earlier comparisons were implicitly against a 1.0 baseline. Stage 4 reports
masked determinism only (Table~\ref{tab:s4-ssh}).

\subsection{API proxy}

A local reverse proxy (\texttt{harness/api\_proxy.py}) sits in front of the
Fireworks endpoint. It logs only \texttt{ts}, \texttt{method},
\texttt{status}, \texttt{duration\_ms}, \texttt{request\_sha256} and
\texttt{usage}; it never logs the \texttt{Authorization} header, the API key,
or request/response body text. A five-command feasibility check on
AdvancedshelLM logged 12 proxy requests, 10 of which matched a command (5 per
timing mode), and measured a median non-API gap of 0.045\,s (min 0.026, max
0.075) over those 5 commands (\texttt{results/exp/smoke/}). Stage 4 runs every target through
the same proxy (Section~\ref{sec:s4-method}).

\section{Stage 4: the 20-run pass}
\label{sec:s4}

\subsection{Method}
\label{sec:s4-method}

Stage 4 is the main result set. Driver: \texttt{harness/run\_stage4.py};
design and schema in \texttt{docs/STAGE4\_NOTES.md}. The pass is reproduced
with
\begin{quote}\footnotesize\ttfamily
python harness/run\_stage4.py --runs 20 --timing prompt --seed 2026 \\
\hspace*{2em}--quiet-check-runs 2 --out results/exp/stage4 \\
python harness/analyze\_stage4.py --stage4-dir results/exp/stage4
\end{quote}
and the data are in \texttt{results/exp/stage4/}. It differs from the baseline
batch in three deliberate ways:

\begin{itemize}[leftmargin=*]
\item \textbf{Prompt timing} instead of quiet timing, so the 0.5\,s settle
  window is not part of the measurement (Section~\ref{sec:floor}).
\item \textbf{Twenty runs per target} instead of five, which supports
  bootstrap confidence intervals resampled by run.
\item \textbf{A fresh container per slot}: each of the 180 slots runs
  \texttt{docker compose down} followed by \texttt{up -d --force-recreate},
  verified 180/180 in \texttt{results/exp/stage4/driver.log}. The baseline
  batch used one long-lived container per target for all five runs.
\end{itemize}

Additional controls: the target order within a run was shuffled with seed
2026 and recorded per row (\texttt{s4\_order\_seed}); runs 1--2 additionally
ran the four LLM SSH targets in quiet mode as a same-day quiet check, written
to \texttt{raw\_quiet/}; credentials were masked at write time; and every
request passed through the API proxy with one slot marker per mode. All 188
slot keys are \texttt{complete}, with no safety stop
(\texttt{results/exp/stage4/status.json}). The pass took 6.85\,h of wall clock
over 180 slots (median 95\,s per slot,
\texttt{results/exp/stage4/analysis\_review.md}).

Budget and cost: 3,699 proxy requests and \$2.0173 at the corrected
cached-token price model (Section~\ref{sec:s4-cost}).

\subsection{Per-target results}

Table~\ref{tab:s4-ssh} is the main results table. Latency excludes timed-out
commands (censored at the 120\,s runner timeout). CIs are 10,000 bootstrap
resamples of run-level medians, seed 20261106. Extra seconds are run-paired
target-minus-ground-truth medians.

\begin{table}[ht]
\centering
\caption{Stage 4 SSH results: 20 runs, prompt timing, 720 scored commands per
target. Sources: \texttt{results/exp/stage4/latency.csv},
\texttt{determinism.csv}, \texttt{bootstrap\_ci.csv},
\texttt{extra\_vs\_gt.csv}; completed counts from
\texttt{analysis\_review.md}.}
\label{tab:s4-ssh}
\scriptsize
\setlength{\tabcolsep}{4pt}
\begin{tabular}{@{}l r r r r r r r@{}}
\toprule
Target & Timeouts & Median s [95\% CI] & p95 s & Extra vs GT s [95\% CI] & Det.\ masked & Fid.\ masked & Completed \\
\midrule
groundtruth\_ssh & 0 & 0.015 [0.015, 0.015] & 0.031 & --- & 0.992 & 0.991 & 720 \\
cowrie & 0 & 0.031 [0.031, 0.031] & 0.047 & 0.016 [0.016, 0.016] & 0.996 & 0.603 & 720 \\
llmhoneypot & 0 & 0.015 [0.015, 0.015] & 0.031 & 0.000 [0.000, 0.000] & 0.996 & 0.437 & 720 \\
AdvancedshelLM & 0 & 2.253 [2.100, 2.437] & 5.311 & 2.216 [2.085, 2.422] & 0.790 & 0.566 & 720 \\
Beelzebub & 4 & 2.308 [2.165, 2.433] & 5.868 & 2.243 [2.150, 2.418] & 0.594 & 0.433 & 716 \\
DECEIVE & 29 & 2.390 [2.228, 2.567] & 8.637 & 2.392 [2.213, 2.552] & 0.509 & 0.482 & 691 \\
VelLMes & 45 & 2.181 [1.995, 2.252] & 4.996 & 2.088 [1.980, 2.237] & 0.848 & 0.637 & 675 \\
\bottomrule
\end{tabular}
\end{table}

The LLM targets sit at 2.18--2.39\,s median, 0.14--0.21\,s apart, with
overlapping confidence intervals; the data do not rank them. Fidelity is
masked similarity against ground truth over the shared-command intersection
(prefixes capped at 4,000 characters);
the llm-honeypot row is the detector target and its fidelity and determinism
are excluded from interpretation.

\textbf{Clock resolution.} Every statistic below 50\,ms is quantised, not
measured. Over 720 completed commands the ground-truth values pile up at
15\,ms (403), 16\,ms (105), 31\,ms (77) and 30\,ms (34); Cowrie and llm-honeypot
show the same 15/16\,ms and 30--32\,ms steps
(\texttt{results/exp/stage4/analysis\_followup.md}~S4). The differences between
ground truth (0.015\,s), llm-honeypot (0.015\,s) and Cowrie (0.031\,s) are one
or two scheduler ticks and are not interpreted.

\subsection{Output quality and prompt integrity}
\label{sec:s4-quality}

Empty outputs out of 720 rows: AdvancedshelLM 103, Cowrie 160, DECEIVE 100,
ground truth 140, llm-honeypot 120, VelLMes 102, Beelzebub 6. Oversized outputs
(over 20,000 characters): DECEIVE 6 (maximum 142,810 characters) and Beelzebub
2; all other targets 0. The run-time leakage counter over
\texttt{cleaned\_output} reports Beelzebub 13 and DECEIVE 2; counting ``interior
prompt followed by text'' in \texttt{cleaned\_output} gives Beelzebub 38/720 in
prompt mode and 6/72 in quiet mode, DECEIVE 2/720, and all other targets 0. The
same pattern on \texttt{raw\_output} gives Beelzebub 88/720 and DECEIVE 2/720,
because the raw capture keeps the trailing prompt and the next command echo (\texttt{results/exp/stage4/analysis.md},
\texttt{analysis\_followup.md}~S1a).

Prompt-mode integrity: \texttt{prompt\_without\_echo} is false in all 5,476 SSH
rows of the pass (5,180 prompt-mode rows, of which 5,040 are scored, plus 296
quiet rows in \texttt{raw\_quiet/}) and in all 6,276 rows including the 800 HTTP
rows; rows without a final prompt are exactly the timed-out
rows plus one Beelzebub quiet row that ends with a leaked next prompt; and no
harness latency is below its matching API duration
(\texttt{analysis\_followup.md}~S1a, S1d). The published Beelzebub leakage case
(14/180) is therefore not the same measurement as the Stage 4 count
(38/720 in \texttt{cleaned\_output}).

The oversize and leakage totals differ from the baseline batch: DECEIVE 6
oversized here versus 2 there, Beelzebub 2 versus 9. Both are prompt-mode
counts, so the comparison is between different timing modes and run counts,
not a measured change.

\subsection{Drift, order, and batch comparison}

Drift controls, Spearman correlation between run index and per-run median
latency over 20 runs (\texttt{results/exp/stage4/drift.csv}): Cowrie
$\rho=0.256$ ($p=0.414$, inconclusive) and ground truth $\rho=0.704$
($p=0.0012$). The ground-truth value is an artefact of two possible run
medians, 15\,ms and 16\,ms; the correlation is one tick flip and is not
reported as drift.

Order effect, Spearman correlation between position-in-run and per-command
latency: AdvancedshelLM $\rho=0.144$ ($p=0.0004$), DECEIVE $\rho=-0.090$
($p=0.049$), VelLMes $\rho=0.075$ ($p=0.12$), Beelzebub $\rho=0.049$
($p=0.34$). The signs disagree and the magnitudes are small; no usable
position effect is claimed (\texttt{results/exp/stage4/analysis.md}).

Table~\ref{tab:s4-vs-v1} sets the baseline batch beside Stage 4. The two
batches differ in timing mode, run count and container lifetime, so the
differences are differences between batches, not effects.

\begin{table}[ht]
\centering
\caption{Baseline (v1) versus Stage 4, same metrics. Sources:
\texttt{results/summary.md}, \texttt{results/exp/summary\_v2.md},
\texttt{results/exp/stage4/latency.csv},
\texttt{results/exp/stage4/analysis\_review.md}.}
\label{tab:s4-vs-v1}
\scriptsize
\setlength{\tabcolsep}{4pt}
\begin{tabular}{@{}l r r r r r r@{}}
\toprule
 & \multicolumn{3}{c}{Timeout rate per 180 commands} & \multicolumn{2}{c}{Median s} & Det.\ masked \\
\cmidrule(lr){2-4}\cmidrule(lr){5-6}
Target & v1 (/180) & Stage 4 & Stage 4 (/180) & v1 & Stage 4 & v1 $\rightarrow$ Stage 4 \\
\midrule
Cowrie & 0 & 0 & 0.0 & 0.563 & 0.031 & 0.996 $\rightarrow$ 0.996 \\
AdvancedshelLM & 0 & 0 & 0.0 & 1.922 & 2.253 & 0.803 $\rightarrow$ 0.790 \\
VelLMes & 9 & 45 & 11.3 & 1.687 & 2.181 & 0.847 $\rightarrow$ 0.848 \\
DECEIVE & 5 & 29 & 7.3 & 1.984 & 2.390 & 0.504 $\rightarrow$ 0.509 \\
Beelzebub & 16 & 4 & 1.0 & 5.023 & 2.308 & 0.682 $\rightarrow$ 0.594 \\
Ground truth & 0 & 0 & 0.0 & 0.562 & 0.015 & 0.984 $\rightarrow$ 0.992 \\
\bottomrule
\end{tabular}
\end{table}

\subsection{HTTP results at equal sample size}

Stage 4 collected 400 galah requests against 400 ground-truth requests. Because
the baseline batch has only 100, the Stage 4 pool was resampled 1,000 times at
100 requests (seed 20261122) to give comparable rates
(\texttt{results/exp/stage4/analysis\_followup.md}~S5).
Table~\ref{tab:galah} sets the two batches side by side at equal sample size.

\begin{table}[ht]
\centering
\caption{Galah HTTP metrics, Stage 4 pool and equal-size resampling versus the
baseline batch. Sources: \texttt{results/exp/stage4/analysis.md};
\texttt{analysis\_followup.md}~S5; \texttt{results/summary.md}.}
\label{tab:galah}
\scriptsize
\begin{tabularx}{\textwidth}{Y r r r r}
\toprule
Metric & Stage 4, 400 req. & Stage 4 per-100 mean & Stage 4 per-100 95\% range & v1, 100 req. \\
\midrule
5xx rate & 0.100 (40/400) & 0.101 & 0.05--0.15 & 0.110 (11/100) \\
Timeout rate & 0.0375 (15/400) & 0.037 & 0.01--0.07 & 0.030 (3/100) \\
Status agreement with nginx & 0.260 (95/366) & 0.260 & 0.187--0.341 & 0.266 (25/94) \\
Distinct \texttt{Server} values & 13 & 6.9 & 4--9 & 6 \\
Repeat \texttt{GET /} similarity & 0.376 & 0.377 & 0.237--0.545 & 0.236 \\
\bottomrule
\end{tabularx}
\end{table}

At equal sampling, the 5xx rate, timeout rate, status agreement and
distinct-\texttt{Server} count of the baseline batch fall inside the Stage 4
per-100 range; repeat similarity (0.236) sits at the lower edge (range starts
at 0.237). The full-400 distinct-\texttt{Server} count of 13 grows with sample
size; the comparable figure is about 7 per 100.

\subsection{Quiet-mode check (runs 1--2)}
\label{sec:quiet}

Runs 1--2 also ran each LLM SSH target in quiet mode, giving 72 prompt/quiet
command pairs per target (\texttt{results/exp/stage4/quiet\_check\_full.csv}). A pair
is flagged when the prompt-mode output is shorter than the quiet-mode output
and its masked similarity is below the Stage 2 quiet-vs-quiet similarity for
that same command.

\begin{table}[ht]
\centering
\caption{Prompt-versus-quiet pairs beyond the Stage 2 noise floor, 72 pairs per
target, 288 in total. Source:
\texttt{results/exp/stage4/quiet\_check\_full.csv}; the analyzer's flagged
subset is \texttt{quiet\_check.csv} (85 rows).}
\label{tab:quiet}
\scriptsize
\begin{tabular}{l r r r}
\toprule
Target & Pairs checked & Beyond the floor & Largest single case \\
\midrule
DECEIVE        & 72 & 11 & \texttt{env}: 684 vs 2,054 chars, similarity 0.28 \\
Beelzebub      & 72 & 6  & \texttt{find\_etc\_conf}: 192 vs 164,049 chars \\
VelLMes        & 72 & 7  & \texttt{ls\_etc} run 2: 0 vs 2,720 chars (timeout) \\
AdvancedshelLM & 72 & 9  & \texttt{passwd}: 83 vs 1,208 chars, similarity 0.12 \\
\bottomrule
\end{tabular}
\end{table}

Table~\ref{tab:quiet} counts them: 33 of the 288 pairs are shorter in prompt
mode and no more similar than their quiet-vs-quiet floor. The full 288-row
table is \texttt{quiet\_check\_full.csv}; the analyzer writes only the pairs it
flags, so \texttt{quiet\_check.csv} holds 85 rows and is not a 288-pair
denominator. Why prompt mode returns less for those commands is UNVERIFIED: the
pairs are different sessions, so output variance and any timing effect are not
separable here.

The VelLMes \texttt{ls\_etc} case is explained by the proxy log: prompt mode
recorded a 120.013\,s harness timeout with 10 bytes captured, while the quiet
run of the same command returned 2,720 characters in 6.0\,s. The proxy row for
that command shows \texttt{status=200}, \texttt{duration\_ms=7203},
\texttt{completion\_tokens=800}: the model produced a capped reply in 7.2\,s and
the target never wrote it to the session
(\texttt{analysis\_followup.md}~S1b). Whether the reply was held or dropped
inside the target is UNVERIFIED; no container-side log was recorded for that
slot.

\subsection{Timeout attribution}
\label{sec:timeouts}

The proxy stamps \texttt{ts} before forwarding and writes the log line after
the response, so file order is completion order. Joining by request start
instead of file order, all 86 timed-out SSH rows of the pass (raw/ and
raw\_quiet/) fall into three classes (Table~\ref{tab:timeouts},
source \texttt{results/exp/stage4/analysis\_timeouts.md}):

\begin{table}[ht]
\centering
\caption{Stage 4 timeout attribution: 86 timed-out rows joined to the proxy log
by request start, each proxy entry used once. Source:
\texttt{results/exp/stage4/analysis\_timeouts.md}.}
\label{tab:timeouts}
\scriptsize
\begin{tabularx}{\textwidth}{l r Y r}
\toprule
Class & n & Description & API duration (s) \\
\midrule
API-side & 9 & request ran past the 120\,s harness timeout & 241.6--291.4, p50 264.1 \\
Target-side & 66 & API answered in time, harness never saw the reply & 1.5--28.1, p50 5.4 \\
Unmatched & 11 & no proxy request assigned to the command & --- \\
\bottomrule
\end{tabularx}
\end{table}

By target: Beelzebub 3 API-side and 3 unmatched; DECEIVE 6 API-side, 15
target-side and 8 unmatched; VelLMes 51 target-side. All nine API-side rows
end at exactly 65,536 completion tokens. The 51 VelLMes rows are the capped
800-token replies: the API answered in 4.5--7.8\,s and the harness saw nothing,
which is the same signature as the \texttt{ls\_etc} case above. Nine of the 11
unmatched rows were commanded while a request of 120\,s or longer from the same
run was still in flight; the other two are DECEIVE run~20 \texttt{whoami} and
\texttt{id}, where no request was logged and the next command took 20.9\,s (a
cold start, not an API timeout). A candidate explanation for the target-side
rows is that the reply is truncated at the cap or dropped before reaching the
terminal; that is UNVERIFIED.

\subsection{max\_tokens enforcement}
\label{sec:maxtokens}

Table~\ref{tab:maxtokens} reports the observed completion-token distribution
per target over the requests attributed to it in Stage 4, with the documented
cap.

\begin{table}[ht]
\centering
\caption{Observed completion tokens per target, Stage 4, versus the documented
cap. Sources: \texttt{results/exp/stage4/analysis\_timeouts.md} (Part A.2);
caps from \texttt{BENCHMARK.md}.}
\label{tab:maxtokens}
\scriptsize
\begin{tabular}{l r r r r r r}
\toprule
Target & Requests & p50 & p95 & Max & Documented cap & Above cap \\
\midrule
AdvancedshelLM & 792 & 51 & 533 & 4,096 & 4,096 & 0 \\
VelLMes & 814 & 56.5 & 800 & 800 & 800 & 0 \\
DECEIVE & 806 & 73 & 1,221 & 65,536 & none & n/a \\
Beelzebub & 789 & 55 & 682 & 65,536 & unknown & n/a \\
galah & 400 & 388 & 3,232 & 65,536 & unknown & n/a \\
\bottomrule
\end{tabular}
\end{table}

Both documented caps hold exactly: AdvancedshelLM peaks at 4,096 and VelLMes at
800, with no request above either. AdvancedshelLM sets the cap in
\texttt{openai\_client.py} line 155 (\texttt{int(os.environ.get('MAX\_TOKENS',
'4096'))}), and every worker call passes it; the two paths that do not are the
Responses-API branch, reachable only for codex-style model names, and a
second provider client with a hard-coded 16384 that the wrapper never selects
(\texttt{docker/advshellm\_ssh\_wrapper.py} passes \texttt{--provider openai}).

The three targets that reach 65,536 are the uncapped ones. An earlier reading
of the proxy log attributed some 65,536-token calls to AdvancedshelLM and
VelLMes; with request-start attribution those calls belong to galah (8),
DECEIVE (6) and Beelzebub (3), so the caps are consistent with the data and no
different call type is required to explain the ceiling
(\texttt{analysis\_timeouts.md}, Part A.2--A.3). Galah's ceiling replies report
\texttt{reasoning\_tokens = 65,536}; DECEIVE's and Beelzebub's spend the ceiling
on visible text. \texttt{BENCHMARK.md} and the v1 configuration table state
4,096 and 800; neither statement is contradicted by the data.

\subsection{API proxy joins and cost}
\label{sec:s4-cost}

All Stage 4 command rows matched a proxy request: 720/720 for each LLM SSH
target and 400/400 for galah, with no unmatched entries. Median API duration was
2.24--2.38\,s for the SSH targets (p95 5.5--9.1\,s) and 3.71\,s for galah (p95
18.2\,s); the median non-API gap was 35--161\,ms for SSH and 46\,ms for galah
(\texttt{results/exp/stage4/proxy\_join.csv}).

Cost with the corrected cached-token model (\$0.15/M uncached prompt, \$0.03/M
cached prompt, \$0.50/M completion) is \$2.0173 over 3,699 requests: 20,459,956
prompt tokens of which 83.72\% were cache hits, and 2,007,696 completion
tokens. Per target: DECEIVE \$0.8983, Beelzebub \$0.4467, VelLMes \$0.2947,
AdvancedshelLM \$0.2253, galah \$0.1522; the four targets that never call the
model cost nothing. The per-run median is \$0.0710, and the two runs above
three times the median (run~4 \$0.2407, run~15 \$0.2268) made the same number of
requests as every other prompt-only run but contained 65,536-token generations
(\texttt{analysis\_followup.md}~S6). An earlier batch, whose driver replayed
run~n n times, was discarded before this pass and is documented in
\texttt{docs/STAGE4\_NOTES.md}~S11; its numbers are not used anywhere here.

\subsection{The Beelzebub observation and the container-lifetime control}
\label{sec:lifetime}

Beelzebub is the one target whose baseline and Stage 4 numbers point in
different directions: 16/180 timeouts and a median of 5.02\,s in the baseline
batch, against 4/720 and 2.31\,s in Stage 4. In the baseline batch its per-run
median rose monotonically across the five runs, 2.16 to 7.09\,s; in Stage 4 the
20 run medians stay between 1.71 and 3.83\,s with no trend
(\texttt{results/exp/stage4/analysis\_followup.md}~S3). The Beelzebub YAML
configuration is identical between the \texttt{final-batch-v1} tag and HEAD
(\texttt{git diff final-batch-v1 HEAD} over the three Beelzebub files is empty);
whether the built binary was identical is not recorded, so that is UNVERIFIED.

The three protocol differences listed in Section~\ref{sec:s4-method} all
changed at once, so the batch difference cannot be assigned to any one of them
from these data. To isolate container lifetime, a control experiment kept one
Beelzebub container for ten consecutive sessions
(\texttt{harness/run\_lifetime\_control.py}, output under
\texttt{results/exp/lifetime\_control/beelzebub\_persistent/}).
Table~\ref{tab:lifetime} summarises it.

\begin{table}[ht]
\centering
\caption{Beelzebub container-lifetime control: ten sessions, one container
(no recreate; container id and \texttt{StartedAt} identical before and after
every run). Source:
\texttt{results/exp/lifetime\_control/beelzebub\_persistent/lifetime.json}.}
\label{tab:lifetime}
\scriptsize
\begin{tabular}{@{}l p{11.2cm}@{}}
\toprule
Item & Value \\
\midrule
Container & one id, \texttt{State.StartedAt} unchanged; \texttt{never\_recreated = true} \\
Sessions & 10, prompt timing, 36 scored commands each \\
Median latency per run & 3.88, 3.95, 4.09, 3.83, 4.13, 5.12, 4.33, 5.43, 5.44, 5.95\,s \\
Timeouts & 0,0,0,0,0,0,2,0,2,2 (6/360), all in runs 7, 9, 10 \\
Spearman, median vs run & $\rho = 0.915$, $p = 0.0006$ (20,000 permutations, seed 20261122) \\
Container memory / log & 11.14 $\rightarrow$ 20.35\,MiB; 38,143 $\rightarrow$ 906,557 bytes \\
Proxy calls / cost & 360 / \$0.8210 \\
\bottomrule
\end{tabular}
\end{table}

The ten-session series is not flat: the median rises from 3.83--3.88\,s to
5.95\,s, and the three sessions that produced a 65,536-token generation (315,
298 and 287\,s) are exactly the three with timeouts. The Stage 4 fresh-container
series for the same target has a median-of-medians of 2.26\,s over 20 runs
(\texttt{analysis\_timeouts.md}, Part B). The persistent series therefore shows
growth and more timeouts than the fresh series, so the hypothesis that a
long-lived container is harmless for this target is not supported. The cause of
the growth is UNVERIFIED: the control differs from Stage 4 in time of day and
machine load, contains a higher share of ceiling generations (3 of 10 sessions
versus 2 of 20 runs), and has only ten sessions.

\section{Findings}

The following statements are limited to what the data above show.

\begin{itemize}[leftmargin=*]
\item Measured masked determinism separated the static baseline (Cowrie 0.996)
  from the LLM SSH targets (0.509--0.848) in Stage 4; ground truth measured
  0.992 (\texttt{results/exp/stage4/determinism.csv}).
\item In Stage 4 the LLM SSH medians were 2.18--2.39\,s against 0.015\,s for
  ground truth, i.e. 2.09--2.39\,s of extra seconds over the run-paired ground
  truth median; values below 50\,ms are clock-quantised and are not
  interpreted (\texttt{extra\_vs\_gt.csv}, \texttt{analysis\_followup.md}~S4).
\item Timeouts were command-specific and repeated across batches: VelLMes
  \texttt{ps\_aux} and \texttt{env} dominate its 45 Stage 4 timeouts (19/20 and
  17/20 runs), and DECEIVE's 29 are led by \texttt{cpuinfo} and \texttt{ps\_aux}
  (6/20 each) (\texttt{analysis\_review.md}).
\item Timeout attribution splits the 86 Stage 4 timed-out rows into 9 API-side
  (all ending at 65,536 completion tokens), 66 target-side and 11 unmatched;
  VelLMes' 51 target-side rows are 800-token replies produced in 4.5--7.8\,s
  that never reached the harness (\texttt{analysis\_timeouts.md}).
\item Observed completion tokens match the documented caps per target
  (AdvancedshelLM 4,096, VelLMes 800) and reach the 65,536 model ceiling only
  for the uncapped targets (\texttt{analysis\_timeouts.md}, Part A.2).
\item Galah's Stage 4 rates are consistent with its baseline rates at equal
  sample size: 5xx 0.101 versus 0.110, status agreement 0.260 versus 0.266
  (\texttt{analysis\_followup.md}~S5).
\item Prompt-mode integrity held across the pass: zero
  \texttt{prompt\_without\_echo} rows in 5,180 prompt-mode SSH rows (5,476 SSH
  rows with the quiet check; 0 of 6,276 rows including HTTP) and no harness
  latency below its matching API duration (\texttt{analysis\_followup.md}~S1).
\item The single-container control shows a rising Beelzebub median
  (3.88 to 5.95\,s, $\rho=0.915$) and 6/360 timeouts, while the fresh-container
  series stays at a median-of-medians of 2.26\,s; the cause is UNVERIFIED
  (\texttt{results/exp/lifetime\_control/}, \texttt{analysis\_timeouts.md}).
\end{itemize}

\section{Limitations}

The benchmark has two batches with different protocols, one shared model and
API endpoint, scripted inputs, censored timeout latencies, and a
character-similarity proxy for fidelity and determinism. Personas, hostnames
and filesystem fixtures differ between targets, so fidelity values fold persona
differences into the score. Shared API load can affect timeout and latency
patterns. The 26-command fidelity intersection is a restricted subset and
host-specific commands are reference-only. HoneyBee has three generations per
cell and static validity checks.

Timing resolution is a hard limit on part of the result set: every value below
50\,ms is quantised by the 15.6\,ms Windows tick, so ground-truth, Cowrie and
detector latencies, and the ``extra seconds'' for Cowrie, are one or two ticks
and are not interpreted.

The batches are not interchangeable. Stage 4 changed timing mode, run count
and container lifetime together, so differences between the batches
(Beelzebub most visibly) are differences between protocols; the lifetime
control addresses one of the three factors with ten sessions and its own
confounds. The response of galah and of the SSH targets to a repeated identical
command is partly a property of the shared endpoint, and galah's cache means
the proxy log is a lower bound on its upstream calls within a slot
(\texttt{docs/STAGE4\_NOTES.md}~S8).

The overlapping confidence intervals do not support a tool ranking. DECEIVE
state-check failures do not establish its state-tracking capability. The
scripted observations do not establish attacker detectability, HoneyBee output
security, or behaviour beyond this model, endpoint and sample. No causal claim
is made anywhere in this report; where a mechanism is suggested it is marked
UNVERIFIED.

\section{Next steps}

Planned work, none of it done in this report:
\begin{enumerate}[leftmargin=*]
\item \textbf{Local-model feasibility check.} One target, a small local model on
  8\,GB RAM, single sessions. The model changes together with the endpoint, so
  no comparison with the hosted batches is implied.
\item \textbf{Persona-matched fidelity.} A control that separates persona and
  filesystem differences from shell errors, so fidelity stops folding the two
  together.
\item \textbf{Repeat of the container-lifetime control with longer sessions and
  matched time of day}, to separate session count from wall-clock and load
  effects.
\item \textbf{Test the VelLMes truncation/drop hypothesis.} The target-side
  signature (capped reply produced, prompt never surfaced) needs a
  container-side log or a capture at the target to say whether the reply is
  held or dropped.
\item \textbf{Add a collaborator-provided honeypot target} to the same manifest.
\item \textbf{Repeat the pass with a second model} while keeping the command and
  request manifests unchanged.
\item \textbf{Add detectability probes} and a separate attacker evaluation.
\end{enumerate}

\begin{thebibliography}{99}

\bibitem{bitterharvest}
Alexander Vetterl and Richard Clayton. ``Bitter Harvest: Systematically
Fingerprinting Low- and Medium-interaction Honeypots at Internet Scale.'' In
\emph{12th USENIX Workshop on Offensive Technologies (WOOT~'18)}, Baltimore,
MD, August 2018. USENIX Association.
% Title, authors, venue, year and location verified against the USENIX WOOT '18
% page on 2026-11-21. UNVERIFIED: page numbers and DOI, which that page omits.

\bibitem{llmintheshell}
Muris Sladi\'c, Veronica Valeros, Carlos Catania, and Sebastian Garcia. ``LLM in
the Shell: Generative Honeypots.'' In \emph{2024 IEEE European Symposium on
Security and Privacy Workshops (EuroS\&PW)}, pp.~430--435.
doi:10.1109/EuroSPW61312.2024.00054.

\bibitem{vellmes-paper}
Muris Sladi\'c, Veronica Valeros, Carlos Catania, and Sebastian Garcia.
``VelLMes: A High-Interaction AI-Based Deception Framework.'' In \emph{2025
IEEE European Symposium on Security and Privacy Workshops (EuroS\&PW)}, pp.~671--679.
doi:10.1109/EuroSPW67616.2025.00082.

\bibitem{honeygpt}
Ziyang Wang, Jianzhou You, Haining Wang, Tianwei Yuan, Shichao Lv, Yang Wang,
and Limin Sun. ``HoneyGPT: Breaking the Trilemma in Terminal Honeypots with
Large Language Model.'' arXiv:2406.01882, 2024.
% Title, author list and 2024 submission year verified against the arXiv abstract
% page on 2026-11-21. UNVERIFIED: peer-reviewed venue, cited here as a preprint;
% the arXiv v2 title reads ``... in Honeypots with Large Language Models''.

\bibitem{reworr}
Reworr and Dmitrii Volkov. ``LLM Agent Honeypot: Monitoring AI Hacking Agents in
the Wild.'' arXiv:2410.13919, 2024.

\bibitem{bridgessok}
Robert A. Bridges, Thomas R. Mitchell, Mauricio Mu\~noz, and Ted Henriksson.
``SoK: Honeypots \& LLMs, More Than the Sum of Their Parts?'' In \emph{2026 IEEE
11th European Symposium on Security and Privacy (EuroS\&P)}, pp.~910--928.
doi:10.1109/EuroSP68448.2026.00063.

\bibitem{galah}
galah. \url{https://github.com/0x4D31/galah}.

\bibitem{deceive}
DECEIVE. \url{https://github.com/splunk/DECEIVE}.

\bibitem{cowrie}
Cowrie. \url{https://github.com/cowrie/cowrie}.

\bibitem{beelzebub}
Beelzebub. \url{https://github.com/beelzebub-labs/beelzebub}.

\bibitem{advancedshellm}
shelLM (AdvancedshelLM). \url{https://github.com/stratosphereips/shelLM}.
% UNVERIFIED: AdvancedshelLM has no standalone repository that resolves; the code
% is part of shelLM.

\bibitem{vellmes}
VelLMes. \url{https://github.com/stratosphereips/VelLMes-AI-Deception-Framework}.

\end{thebibliography}

\end{document}
